\documentclass[12pt,a4paper]{article}
\usepackage[utf8]{inputenc}
\usepackage[T1]{fontenc}
\usepackage{amsmath}
\usepackage{amsfonts}
\usepackage{amssymb}
\usepackage{graphicx}
\usepackage[margin=1in]{geometry}
\usepackage{array}
\usepackage{longtable}
\usepackage{booktabs}
\usepackage{listings}
\usepackage{xcolor}
\usepackage{enumitem}
\usepackage{hyperref}
\usepackage{tikz}
\usepackage{titlesec}

\definecolor{codegreen}{rgb}{0,0.6,0}
\definecolor{codegray}{rgb}{0.5,0.5,0.5}
\definecolor{codepurple}{rgb}{0.58,0,0.82}
\definecolor{backcolour}{rgb}{0.95,0.95,0.92}

\lstdefinestyle{mystyle}{
    backgroundcolor=\color{backcolour},
    commentstyle=\color{codegreen},
    keywordstyle=\color{magenta},
    numberstyle=\tiny\color{codegray},
    stringstyle=\color{codepurple},
    basicstyle=\ttfamily\footnotesize,
    breakatwhitespace=false,
    breaklines=true,
    captionpos=b,
    keepspaces=true,
    numbers=left,
    numbersep=5pt,
    showspaces=false,
    showstringspaces=false,
    showtabs=false,
    tabsize=2
}

\lstset{style=mystyle}

\titleformat{\section}{\large\bfseries}{\thesection}{1em}{}
\titleformat{\subsection}{\normalsize\bfseries}{\thesubsection}{1em}{}
\titleformat{\subsubsection}{\normalsize\bfseries}{\thesubsubsection}{1em}{}

\title{Module 1: Crystal Physics \& Reciprocal Lattice\\ \large Complete Worked Solutions \& TA Rubric}
\author{PH7302E Condensed Matter Physics \\ Department of Physics \\ National Institute of Technology Calicut}
\date{Academic Year 2026--27}

\begin{document}

\maketitle

\tableofcontents
\newpage

\section*{Preface for Teaching Assistants}
This document serves as a comprehensive solution rubric and grading guide for Module 1 of PH7302E Condensed Matter Physics. It contains:

\begin{enumerate}
    \item Complete solutions to all assignment problems
    \item Detailed grading rubrics with point allocations
    \item Common student errors and misconceptions
    \item TA guidelines for consistent grading
    \item Supplementary teaching materials
\end{enumerate}

\textbf{Total Points for Assignment:} 100

\section{Grading Policy}
\begin{itemize}
    \item Each problem is worth the points indicated
    \item Partial credit should be awarded for:
    \begin{itemize}
        \item Correct approach even if final answer is wrong
        \item Correct setup of equations
        \item Appropriate use of formulas
        \item Clear explanation of reasoning
    \end{itemize}
    \item Deductions should be made for:
    \begin{itemize}
        \item Missing units
        \item Incorrect significant figures
        \item Lack of clarity in derivations
        \item Missing steps in proofs
    \end{itemize}
\end{itemize}

\newpage

\section{Assignment 1: Complete Solutions}

\subsection*{Lecture 1 Problems}

\subsubsection*{Problem 1: Packing Fraction of FCC (5 points)}
\textbf{Solution:}
In FCC, atoms touch along face diagonals. If $a$ is the lattice constant and $r$ is the atomic radius:
\[4r = a\sqrt{2} \implies r = \frac{a\sqrt{2}}{4}\]

Volume of atoms per unit cell:
\begin{itemize}
    \item 4 atoms per unit cell
    \item Volume of one atom = $\frac{4}{3}\pi r^3$
    \item Total atomic volume = $4 \times \frac{4}{3}\pi r^3 = \frac{16}{3}\pi r^3$
\end{itemize}

Packing fraction:
\[\eta = \frac{4 \times \frac{4}{3}\pi r^3}{a^3} = \frac{16\pi r^3}{3a^3} = \frac{16\pi}{3} \left(\frac{a\sqrt{2}}{4}\right)^3 \frac{1}{a^3}\]
\[\eta = \frac{16\pi}{3} \times \frac{2\sqrt{2}}{64} = \frac{32\sqrt{2}\pi}{192} = \frac{\pi}{3\sqrt{2}} \approx 0.7405\]

\textbf{Grading Rubric:}
\begin{table}[h]
\centering
\begin{tabular}{|p{10cm}|c|}
\hline
\textbf{Criterion} & \textbf{Points} \\
\hline
Correct relationship between $r$ and $a$ ($4r = a\sqrt{2}$) & 2 \\
Correct number of atoms per unit cell (4) & 1 \\
Correct formula for packing fraction & 1 \\
Correct final result ($\pi/3\sqrt{2} \approx 0.74$) & 1 \\
\hline
\end{tabular}
\end{table}

\textbf{Common Errors:}
\begin{itemize}
    \item Using incorrect relationship between $r$ and $a$
    \item Forgetting to multiply by number of atoms
    \item Using wrong number of atoms per unit cell
\end{itemize}

\subsubsection*{Problem 2: Atoms in Diamond Conventional Cell (3 points)}
\textbf{Solution:}
Diamond structure: FCC lattice with 2-atom basis at (0,0,0) and $(\frac{1}{4}, \frac{1}{4}, \frac{1}{4})$.

FCC has 4 lattice points per conventional cell.
Each lattice point has 2 atoms associated with it.
Total atoms = $4 \times 2 = 8$ atoms per conventional cell.

\textbf{Grading Rubric:}
\begin{table}[h]
\centering
\begin{tabular}{|p{10cm}|c|}
\hline
\textbf{Criterion} & \textbf{Points} \\
\hline
Recognizes FCC with 4 lattice points per cell & 1 \\
Recognizes 2-atom basis & 1 \\
Correct final answer (8 atoms) & 1 \\
\hline
\end{tabular}
\end{table}

\subsubsection*{Problem 3: (110) Plane in BCC Lattice (4 points)}
\textbf{Solution:}
The (110) plane in a BCC lattice contains:
\begin{itemize}
    \item Corner atoms at (0,0,0), (1,0,0), (0,1,0), (1,1,0) in the plane
    \item Body center atoms at ($\frac{1}{2}, \frac{1}{2}, 0$) in the plane
    \item The plane cuts through the body-centered atom
\end{itemize}

\textbf{Grading Rubric:}
\begin{table}[h]
\centering
\begin{tabular}{|p{10cm}|c|}
\hline
\textbf{Criterion} & \textbf{Points} \\
\hline
Correct identification of plane intercepts & 2 \\
Correct identification of atoms in plane & 1 \\
Clear sketch with labeling & 1 \\
\hline
\end{tabular}
\end{table}

\subsubsection*{Problem 4: Miller Indices from Intercepts (4 points)}
\textbf{Solution:}
Intercepts: $2a, 3b, 6c$
\begin{enumerate}
    \item Intercepts in units of lattice constants: $(2, 3, 6)$
    \item Take reciprocals: $(\frac{1}{2}, \frac{1}{3}, \frac{1}{6})$
    \item Clear fractions (multiply by 6): $(3, 2, 1)$
    \item Miller indices: $(321)$
\end{enumerate}

\textbf{Grading Rubric:}
\begin{table}[h]
\centering
\begin{tabular}{|p{10cm}|c|}
\hline
\textbf{Criterion} & \textbf{Points} \\
\hline
Correct intercepts in units of lattice constants & 1 \\
Correct reciprocals & 1 \\
Correct clearing of fractions & 1 \\
Correct final Miller indices & 1 \\
\hline
\end{tabular}
\end{table}

\subsubsection*{Problem 5: Interplanar Spacing for (220) Planes in Copper (4 points)}
\textbf{Solution:}
For cubic lattice:
\[d_{220} = \frac{a}{\sqrt{h^2 + k^2 + l^2}} = \frac{3.61}{\sqrt{2^2 + 2^2 + 0^2}} = \frac{3.61}{\sqrt{8}} = \frac{3.61}{2\sqrt{2}} = 1.28 \text{ \AA}\]

\textbf{Grading Rubric:}
\begin{table}[h]
\centering
\begin{tabular}{|p{10cm}|c|}
\hline
\textbf{Criterion} & \textbf{Points} \\
\hline
Correct formula for cubic interplanar spacing & 1 \\
Correct values for h, k, l & 1 \\
Correct arithmetic & 1 \\
Correct units & 1 \\
\hline
\end{tabular}
\end{table}

\subsubsection*{Problem 6: $(11\bar{2})$ Plane (4 points)}
\textbf{Solution:}
For the $(11\bar{2})$ plane:
\begin{itemize}
    \item Intercepts: $(\frac{1}{1}, \frac{1}{1}, \frac{1}{-2}) = (1, 1, -\frac{1}{2})$
    \item The plane intersects the x-axis at $a$, y-axis at $b$, and z-axis at $-c/2$
    \item The negative intercept means the plane cuts the negative z-axis
\end{itemize}

\textbf{Grading Rubric:}
\begin{table}[h]
\centering
\begin{tabular}{|p{10cm}|c|}
\hline
\textbf{Criterion} & \textbf{Points} \\
\hline
Correct intercepts from Miller indices & 2 \\
Correct drawing with axes labeled & 1 \\
Correct indication of negative intercept & 1 \\
\hline
\end{tabular}
\end{table}

\subsection*{Lecture 2 Problems}

\subsubsection*{Problem 7: Reciprocal Lattice of 2D Hexagonal Lattice (8 points)}
\textbf{Solution:}
Given: $\mathbf{a}_1 = a(1,0)$, $\mathbf{a}_2 = a(\frac{1}{2}, \frac{\sqrt{3}}{2})$

For 2D, using the vector cross product with $\hat{z}$:
\[V_{2D} = |\mathbf{a}_1 \times \mathbf{a}_2| = a^2 \frac{\sqrt{3}}{2}\]

\[\mathbf{b}_1 = 2\pi \frac{\mathbf{a}_2 \times \hat{z}}{V_{2D}} = 2\pi \frac{a(\frac{1}{2}, \frac{\sqrt{3}}{2}, 0) \times (0,0,1)}{a^2\sqrt{3}/2}\]

\[= 2\pi \frac{a(\frac{\sqrt{3}}{2}, -\frac{1}{2}, 0)}{a^2\sqrt{3}/2} = \frac{2\pi}{a}(1, -\frac{1}{\sqrt{3}})\]

\[\mathbf{b}_2 = 2\pi \frac{\hat{z} \times \mathbf{a}_1}{V_{2D}} = 2\pi \frac{(0,0,1) \times a(1,0,0)}{a^2\sqrt{3}/2} = \frac{2\pi}{a}(0, \frac{2}{\sqrt{3}})\]

The angle between $\mathbf{b}_1$ and $\mathbf{b}_2$ is 60°:
\[\cos \theta = \frac{\mathbf{b}_1 \cdot \mathbf{b}_2}{|\mathbf{b}_1||\mathbf{b}_2|} = \frac{2\pi/a \times 0 + (-2\pi/(a\sqrt{3})) \times (4\pi/(a\sqrt{3}))}{(2\pi/a)(2\pi/a)(2/\sqrt{3})} = \frac{1}{2}\]
Therefore $\theta = 60^\circ$, confirming hexagonal symmetry.

\textbf{Grading Rubric:}
\begin{table}[h]
\centering
\begin{tabular}{|p{10cm}|c|}
\hline
\textbf{Criterion} & \textbf{Points} \\
\hline
Correct 2D reciprocal lattice definition & 2 \\
Correct calculation of $V_{2D}$ & 1 \\
Correct calculation of $\mathbf{b}_1$ & 1.5 \\
Correct calculation of $\mathbf{b}_2$ & 1.5 \\
Correct angle calculation & 1 \\
Conclusion: reciprocal lattice is hexagonal & 1 \\
\hline
\end{tabular}
\end{table}

\subsubsection*{Problem 8: Reciprocal of BCC is FCC (8 points)}
\textbf{Solution:}
Given BCC direct primitive vectors:
\[\mathbf{a}_1 = \frac{a}{2}(-\hat{x} + \hat{y} + \hat{z}), \quad \mathbf{a}_2 = \frac{a}{2}(\hat{x} - \hat{y} + \hat{z}), \quad \mathbf{a}_3 = \frac{a}{2}(\hat{x} + \hat{y} - \hat{z})\]

\begin{align*}
\mathbf{a}_2 \times \mathbf{a}_3 &= \frac{a^2}{4}
\begin{vmatrix}
\hat{x} & \hat{y} & \hat{z} \\
1 & -1 & 1 \\
1 & 1 & -1
\end{vmatrix} \\
&= \frac{a^2}{4}[\hat{x}((-1)(-1) - (1)(1)) - \hat{y}((1)(-1) - (1)(1)) + \hat{z}((1)(1) - (-1)(1))] \\
&= \frac{a^2}{4}[\hat{x}(1-1) - \hat{y}(-1-1) + \hat{z}(1+1)] \\
&= \frac{a^2}{4}[0\hat{x} + 2\hat{y} + 2\hat{z}] = \frac{a^2}{2}(\hat{y} + \hat{z})
\end{align*}

$V = \mathbf{a}_1 \cdot (\mathbf{a}_2 \times \mathbf{a}_3) = \frac{a}{2}(-\hat{x} + \hat{y} + \hat{z}) \cdot \frac{a^2}{2}(\hat{y} + \hat{z}) = \frac{a^3}{4}(1+1) = \frac{a^3}{2}$

\[\mathbf{b}_1 = 2\pi \frac{\mathbf{a}_2 \times \mathbf{a}_3}{V} = 2\pi \frac{a^2(\hat{y} + \hat{z})/2}{a^3/2} = \frac{2\pi}{a}(\hat{y} + \hat{z})\]

Similarly:
\[\mathbf{b}_2 = \frac{2\pi}{a}(\hat{x} + \hat{z}), \quad \mathbf{b}_3 = \frac{2\pi}{a}(\hat{x} + \hat{y})\]

These are FCC primitive vectors with lattice constant $4\pi/a$, proving the reciprocal of BCC is FCC.

\textbf{Grading Rubric:}
\begin{table}[h]
\centering
\begin{tabular}{|p{10cm}|c|}
\hline
\textbf{Criterion} & \textbf{Points} \\
\hline
Correct BCC primitive vectors & 1 \\
Correct calculation of cross products & 2 \\
Correct calculation of volume $V$ & 1 \\
Correct calculation of $\mathbf{b}_1, \mathbf{b}_2, \mathbf{b}_3$ & 2 \\
Recognition of FCC structure & 1 \\
Correct lattice constant ($4\pi/a$) & 1 \\
\hline
\end{tabular}
\end{table}

\subsubsection*{Problem 9: $\mathbf{G}_{hkl}$ Perpendicular to $(hkl)$ Plane (6 points)}
\textbf{Solution:}
Given: $\mathbf{G}_{hkl} = h\mathbf{b}_1 + k\mathbf{b}_2 + l\mathbf{b}_3$

The $(hkl)$ plane contains vectors such as $\mathbf{a}_1/h - \mathbf{a}_2/k$ (both endpoints lie in the plane).

\begin{align*}
\mathbf{G}_{hkl} \cdot (\mathbf{a}_1/h - \mathbf{a}_2/k) &= (h\mathbf{b}_1 + k\mathbf{b}_2 + l\mathbf{b}_3) \cdot (\mathbf{a}_1/h - \mathbf{a}_2/k) \\
&= h\mathbf{b}_1 \cdot \mathbf{a}_1/h - h\mathbf{b}_1 \cdot \mathbf{a}_2/k + k\mathbf{b}_2 \cdot \mathbf{a}_1/h - k\mathbf{b}_2 \cdot \mathbf{a}_2/k \\
&+ l\mathbf{b}_3 \cdot \mathbf{a}_1/h - l\mathbf{b}_3 \cdot \mathbf{a}_2/k
\end{align*}

Using $\mathbf{a}_i \cdot \mathbf{b}_j = 2\pi \delta_{ij}$:
\[= 2\pi - 0 + 0 - 2\pi + 0 - 0 = 0\]

Therefore $\mathbf{G}_{hkl}$ is perpendicular to the $(hkl)$ plane.

\textbf{Grading Rubric:}
\begin{table}[h]
\centering
\begin{tabular}{|p{10cm}|c|}
\hline
\textbf{Criterion} & \textbf{Points} \\
\hline
Correct expression for $\mathbf{G}_{hkl}$ & 1 \\
Correct identification of vectors in the plane & 1.5 \\
Correct application of orthogonality relation & 1.5 \\
Correct final conclusion ($\mathbf{G} \perp$ plane) & 2 \\
\hline
\end{tabular}
\end{table}

\subsubsection*{Problem 10: Bragg Angle for NaCl (200) (5 points)}
\textbf{Solution:}
For cubic: $d_{200} = \frac{a}{\sqrt{2^2 + 0^2 + 0^2}} = \frac{5.64}{2} = 2.82 \text{ \AA}$

Using Bragg's law: $2d\sin\theta = n\lambda$
For first order ($n=1$):
\[\sin\theta = \frac{\lambda}{2d} = \frac{1.54}{2 \times 2.82} = \frac{1.54}{5.64} = 0.2730\]

\[\theta = \sin^{-1}(0.2730) = 15.84^\circ\]
\[2\theta = 31.68^\circ\]

\textbf{Grading Rubric:}
\begin{table}[h]
\centering
\begin{tabular}{|p{10cm}|c|}
\hline
\textbf{Criterion} & \textbf{Points} \\
\hline
Correct $d_{200}$ calculation & 1.5 \\
Correct application of Bragg's law & 1.5 \\
Correct arithmetic & 1 \\
Correct final angles with units & 1 \\
\hline
\end{tabular}
\end{table}

\subsubsection*{Problem 11: Structure Factor of Diamond (6 points)}
\textbf{Solution:}
Diamond structure: FCC basis with second atom at $(\frac{1}{4}, \frac{1}{4}, \frac{1}{4})$

Structure factor:
\[F_{hkl} = f_{Si}[1 + e^{2\pi i(h/4 + k/4 + l/4)}] \times F_{FCC}\]

where $F_{FCC} = 1 + e^{\pi i(k+l)} + e^{\pi i(h+l)} + e^{\pi i(h+k)}$

$F_{FCC}$ is zero when $h,k,l$ are mixed.

When $h,k,l$ are all even or all odd:
\[F_{hkl} = f_{Si}[1 + e^{2\pi i(h+k+l)/4}]F_{FCC}\]

For $h + k + l = 4n + 2$:
\[e^{2\pi i(4n+2)/4} = e^{\pi i(2n+1)} = -1\]

Therefore $1 + e^{2\pi i(h+k+l)/4} = 1 - 1 = 0$

Hence $F_{hkl} = 0$ when $h + k + l = 4n + 2$.

\textbf{Grading Rubric:}
\begin{table}[h]
\centering
\begin{tabular}{|p{10cm}|c|}
\hline
\textbf{Criterion} & \textbf{Points} \\
\hline
Correct diamond structure basis & 1 \\
Correct expression for $F_{hkl}$ & 1.5 \\
Correct $F_{FCC}$ expression & 1 \\
Correct analysis of $F_{FCC}$ zeros & 0.5 \\
Correct analysis of second atom contribution & 1 \\
Correct final condition ($h+k+l = 4n+2$) & 1 \\
\hline
\end{tabular}
\end{table}

\subsubsection*{Problem 12: (100) Absence in BCC (4 points)}
\textbf{Solution:}
For BCC: basis atoms at (0,0,0) and $(\frac{1}{2}, \frac{1}{2}, \frac{1}{2})$

Structure factor:
\[F_{100} = f[1 + e^{\pi i(1+0+0)}] = f[1 + e^{\pi i}] = f[1 - 1] = 0\]

The (100) reflection is absent because the waves scattered from the corner atoms and the body center atom are out of phase by $\pi$, causing destructive interference.

For SC with one atom basis at (0,0,0):
\[F_{100} = f \neq 0\]

Thus the (100) reflection is present in SC but absent in BCC.

\textbf{Grading Rubric:}
\begin{table}[h]
\centering
\begin{tabular}{|p{10cm}|c|}
\hline
\textbf{Criterion} & \textbf{Points} \\
\hline
Correct BCC structure factor expression & 1.5 \\
Correct evaluation of $F_{100} = 0$ & 1 \\
Correct SC structure factor evaluation & 0.5 \\
Physical explanation & 1 \\
\hline
\end{tabular}
\end{table}

\subsection*{Lecture 3 Problems}

\subsubsection*{Problem 13: Powder Pattern Indexing for Iron (6 points)}
\textbf{Solution:}
For BCC, allowed reflections require $h+k+l$ even.

Using $\sin^2\theta = \frac{\lambda^2}{4a^2}(h^2 + k^2 + l^2)$

For BCC, the first allowed reflections:
\begin{center}
\begin{tabular}{|c|c|c|}
\hline
$(hkl)$ & $h^2+k^2+l^2$ & $\sin^2\theta$ \\
\hline
(110) & 2 & $\frac{\lambda^2}{4a^2} \times 2$ \\
(200) & 4 & $\frac{\lambda^2}{4a^2} \times 4$ \\
(211) & 6 & $\frac{\lambda^2}{4a^2} \times 6$ \\
(220) & 8 & $\frac{\lambda^2}{4a^2} \times 8$ \\
(310) & 10 & $\frac{\lambda^2}{4a^2} \times 10$ \\
(222) & 12 & $\frac{\lambda^2}{4a^2} \times 12$ \\
\hline
\end{tabular}
\end{center}

First five lines: (110), (200), (211), (220), (310)

\textbf{Grading Rubric:}
\begin{table}[h]
\centering
\begin{tabular}{|p{10cm}|c|}
\hline
\textbf{Criterion} & \textbf{Points} \\
\hline
Correct BCC selection rule ($h+k+l$ even) & 1 \\
Correct calculation of $h^2+k^2+l^2$ for each reflection & 2 \\
Correct order of reflections & 2 \\
Correct identification of first five lines & 1 \\
\hline
\end{tabular}
\end{table}

\subsubsection*{Problem 14: Why Powder Method Works (4 points)}
\textbf{Solution:}
The powder method works because:
\begin{enumerate}
    \item The powder contains crystallites in all possible orientations
    \item For any family of planes $(hkl)$, there will always be some crystallites oriented at the correct Bragg angle $\theta$
    \item Each set of planes produces a cone of diffracted radiation
    \item The random orientation ensures that all allowed reflections appear in the pattern
\end{enumerate}

\textbf{Grading Rubric:}
\begin{table}[h]
\centering
\begin{tabular}{|p{10cm}|c|}
\hline
\textbf{Criterion} & \textbf{Points} \\
\hline
Random orientation of crystallites & 1 \\
Some crystallites always satisfy Bragg's law & 1 \\
Cone formation explanation & 1 \\
Complete coverage of all reflections & 1 \\
\hline
\end{tabular}
\end{table}

\subsubsection*{Problem 15: Multiplicity Factor for {110} (4 points)}
\textbf{Solution:}
For {110} planes in cubic crystals, the equivalent planes are obtained by permutations and sign changes of the indices.

Number of permutations of (1,1,0) with two different indices = $\frac{3!}{2!} = 3$

For each permutation, sign changes:
\begin{itemize}
    \item For indices with no zeros: $2^3$ sign combinations possible
    \item But $(hkl)$ and $(\bar{h}\bar{k}\bar{l})$ are the same plane, so divide by 2
\end{itemize}

Multiplicity = $3 \times 2^2 = 12$ (since the zero index has only 1 sign)

Alternative: {110} has 12 equivalent planes.

\textbf{Grading Rubric:}
\begin{table}[h]
\centering
\begin{tabular}{|p{10cm}|c|}
\hline
\textbf{Criterion} & \textbf{Points} \\
\hline
Correct number of permutations & 1.5 \\
Correct treatment of sign changes & 1.5 \\
Correct final multiplicity (12) & 1 \\
\hline
\end{tabular}
\end{table}

\subsubsection*{Problem 16: 5-fold Rotation Incompatibility (5 points)}
\textbf{Solution:}
Consider a lattice row with spacing $a$. A $C_n$ rotation by angle $2\pi/n$ maps lattice points to lattice points.

After rotation, the projection of the rotated vector onto the original row must be an integer multiple of $a$:
\[a\cos(2\pi/n) = ma \implies \cos(2\pi/n) = m\]

where $m$ is an integer.

For $n=5$: $\cos(72^\circ) = 0.3090$, not an integer.

Therefore 5-fold rotation is incompatible with translational symmetry.

The only possible values are:
\begin{itemize}
    \item $n=1$: $\cos(360^\circ) = 1$
    \item $n=2$: $\cos(180^\circ) = -1$
    \item $n=3$: $\cos(120^\circ) = -1/2$
    \item $n=4$: $\cos(90^\circ) = 0$
    \item $n=6$: $\cos(60^\circ) = 1/2$
\end{itemize}

\textbf{Grading Rubric:}
\begin{table}[h]
\centering
\begin{tabular}{|p{10cm}|c|}
\hline
\textbf{Criterion} & \textbf{Points} \\
\hline
Correct setup with lattice spacing $a$ & 1 \\
Correct condition $\cos(2\pi/n)$ must be integer & 1.5 \\
Correct evaluation for $n=5$ & 1 \\
Correct listing of allowed values & 1.5 \\
\hline
\end{tabular}
\end{table}

\subsubsection*{Problem 17: Point Group of Cube (3 points)}
\textbf{Solution:}
A cube has point group $O_h$ (Schoenflies notation) or $m\overline{3}m$ (International notation).

Operations: 48 total
\begin{itemize}
    \item 24 proper rotations (including identity)
    \item 24 improper rotations (including inversion)
\end{itemize}

\textbf{Grading Rubric:}
\begin{table}[h]
\centering
\begin{tabular}{|p{10cm}|c|}
\hline
\textbf{Criterion} & \textbf{Points} \\
\hline
Correct Schoenflies notation ($O_h$) & 1 \\
Correct International notation ($m\overline{3}m$) & 1 \\
Correct number of operations (48) & 1 \\
\hline
\end{tabular}
\end{table}

\subsubsection*{Problem 18: Space Groups in 2D (2 points)}
\textbf{Solution:}
There are 17 space groups in 2D.

\textbf{Grading Rubric:}
\begin{table}[h]
\centering
\begin{tabular}{|p{10cm}|c|}
\hline
\textbf{Criterion} & \textbf{Points} \\
\hline
Correct answer (17) & 2 \\
\hline
\end{tabular}
\end{table}

\subsection*{Lecture 4 Problems}

\subsubsection*{Problem 19: Copper Properties (9 points)}
\textbf{Solution:}
(a) In FCC, atoms touch along face diagonal: $4r = \sqrt{2}a$
\[r = \frac{\sqrt{2}a}{4} = \frac{\sqrt{2} \times 3.61}{4} = 1.28 \text{ \AA}\]

(b) Density:
\[\rho = \frac{4M}{N_A a^3}\]
\[= \frac{4 \times 63.55 \text{ g/mol}}{6.022 \times 10^{23} \text{ mol}^{-1} \times (3.61 \times 10^{-8} \text{ cm})^3}\]
\[= \frac{254.2}{6.022 \times 10^{23} \times 4.70 \times 10^{-23}} = \frac{254.2}{28.3} = 8.96 \text{ g/cm}^3\]

(c) $d_{111} = \frac{a}{\sqrt{3}} = \frac{3.61}{\sqrt{3}} = 2.08 \text{ \AA}$

\textbf{Grading Rubric:}
\begin{table}[h]
\centering
\begin{tabular}{|p{10cm}|c|}
\hline
\textbf{Criterion} & \textbf{Points} \\
\hline
(a) Correct relationship and calculation of $r$ & 3 \\
(b) Correct formula with correct values & 2 \\
(b) Correct arithmetic and units & 1 \\
(c) Correct formula and calculation & 2 \\
(c) Correct units & 1 \\
\hline
\end{tabular}
\end{table}

\subsubsection*{Problem 20: BCC Diffraction (6 points)}
\textbf{Solution:}
Given: $2\theta = 30^\circ$, $\lambda = 1.54$ \AA, (110) planes

$\theta = 15^\circ$

Bragg's law: $2d_{110}\sin\theta = \lambda$
\[d_{110} = \frac{\lambda}{2\sin\theta} = \frac{1.54}{2\sin(15^\circ)} = \frac{1.54}{2 \times 0.2588} = \frac{1.54}{0.5176} = 2.98 \text{ \AA}\]

For cubic: $d_{110} = \frac{a}{\sqrt{1^2 + 1^2 + 0^2}} = \frac{a}{\sqrt{2}}$

Therefore: $a = d_{110}\sqrt{2} = 2.98 \times 1.414 = 4.21 \text{ \AA}$

\textbf{Grading Rubric:}
\begin{table}[h]
\centering
\begin{tabular}{|p{10cm}|c|}
\hline
\textbf{Criterion} & \textbf{Points} \\
\hline
Correct Bragg's law application & 2 \\
Correct calculation of $d_{110}$ & 2 \\
Correct relationship between $a$ and $d_{110}$ & 1 \\
Correct final value with units & 1 \\
\hline
\end{tabular}
\end{table}

\subsubsection*{Problem 21: 2D Hexagonal Lattice (6 points)}
\textbf{Solution:}
Given: $\mathbf{a}_1 = a(1,0)$, $\mathbf{a}_2 = a(\frac{1}{2}, \frac{\sqrt{3}}{2})$

The reciprocal lattice vectors have already been derived in Problem 7:
\[\mathbf{b}_1 = \frac{2\pi}{a}(1, -\frac{1}{\sqrt{3}}), \quad \mathbf{b}_2 = \frac{2\pi}{a}(0, \frac{2}{\sqrt{3}})\]

Lengths:
\[|\mathbf{b}_1| = \frac{2\pi}{a}\sqrt{1 + \frac{1}{3}} = \frac{2\pi}{a}\frac{2}{\sqrt{3}}\]
\[|\mathbf{b}_2| = \frac{2\pi}{a}\frac{2}{\sqrt{3}}\]

The angle between $\mathbf{b}_1$ and $\mathbf{b}_2$:
\[\cos\theta = \frac{\mathbf{b}_1 \cdot \mathbf{b}_2}{|\mathbf{b}_1||\mathbf{b}_2|} = \frac{(\frac{2\pi}{a})(-\frac{1}{\sqrt{3}}) \times (\frac{2\pi}{a})(\frac{2}{\sqrt{3}})}{(\frac{2\pi}{a}\frac{2}{\sqrt{3}})^2} = \frac{1}{2}\]

Therefore $\theta = 60^\circ$, confirming hexagonal symmetry.

\textbf{Grading Rubric:}
\begin{table}[h]
\centering
\begin{tabular}{|p{10cm}|c|}
\hline
\textbf{Criterion} & \textbf{Points} \\
\hline
Correct $\mathbf{b}_1$ and $\mathbf{b}_2$ & 2 \\
Correct lengths calculation & 1 \\
Correct angle calculation & 2 \\
Conclusion: reciprocal lattice is hexagonal & 1 \\
\hline
\end{tabular}
\end{table}

\subsubsection*{Problem 22: Sodium BCC (6 points)}
\textbf{Solution:}
Given: $2\theta = 29.8^\circ$, $\lambda = 1.54$ \AA

$\theta = 14.9^\circ$

Bragg's law: $2d\sin\theta = \lambda$
\[d = \frac{\lambda}{2\sin\theta} = \frac{1.54}{2\sin(14.9^\circ)} = \frac{1.54}{2 \times 0.2571} = \frac{1.54}{0.5142} = 2.99 \text{ \AA}\]

For BCC, first reflection is from (110) planes.
\[d_{110} = \frac{a}{\sqrt{2}}\]

Therefore: $a = d_{110}\sqrt{2} = 2.99 \times 1.414 = 4.23 \text{ \AA}$

The reflecting planes are (110).

\textbf{Grading Rubric:}
\begin{table}[h]
\centering
\begin{tabular}{|p{10cm}|c|}
\hline
\textbf{Criterion} & \textbf{Points} \\
\hline
Correct Bragg's law application & 2 \\
Correct calculation of $d$ & 1.5 \\
Correct identification of BCC first reflection & 1 \\
Correct $a$ calculation & 1 \\
Correct identification of planes & 0.5 \\
\hline
\end{tabular}
\end{table}

\subsubsection*{Problem 23: Zinc Blende Structure Factor (7 points)}
\textbf{Solution:}
Zinc blende: FCC lattice with basis of two different atoms at (0,0,0) and $(\frac{1}{4}, \frac{1}{4}, \frac{1}{4})$.

Structure factor:
\[F_{hkl} = f_{Zn}F_{FCC} + f_S e^{2\pi i(h+k+l)/4}F_{FCC}\]

where $F_{FCC} = 1 + e^{\pi i(k+l)} + e^{\pi i(h+l)} + e^{\pi i(h+k)}$

\[F_{hkl} = F_{FCC}[f_{Zn} + f_S e^{2\pi i(h+k+l)/4}]\]

Systematic absences:
1. When $h,k,l$ are mixed (FCC structure factor zero): $F_{hkl} = 0$
2. When $h,k,l$ are all even or all odd:
   \[F_{hkl} = F_{FCC}[f_{Zn} + f_S e^{2\pi i(h+k+l)/4}]\]
   Additional absences occur when $f_{Zn} + f_S e^{2\pi i(h+k+l)/4} = 0$

\textbf{Grading Rubric:}
\begin{table}[h]
\centering
\begin{tabular}{|p{10cm}|c|}
\hline
\textbf{Criterion} & \textbf{Points} \\
\hline
Correct zinc blende basis & 1 \\
Correct $F_{FCC}$ expression & 1.5 \\
Correct combined structure factor expression & 2 \\
Correct mixed indices condition & 1 \\
Correct additional absence condition & 1.5 \\
\hline
\end{tabular}
\end{table}

\subsubsection*{Problem 24: Neumann's Principle Application (5 points)}
\textbf{Solution:}
Neumann's principle states that the symmetry elements of any physical property must include the symmetry elements of the point group of the crystal.

For a cubic crystal with inversion symmetry (e.g., NaCl):
\begin{itemize}
    \item Inversion symmetry requires any physical property to be invariant under $\mathbf{r} \rightarrow -\mathbf{r}$
    \item Piezoelectricity requires a change in polarization under strain: $P_i = d_{ijk}\sigma_{jk}$
    \item Under inversion, $P \rightarrow -P$ and $\sigma \rightarrow \sigma$ (stress is a polar tensor)
    \item Therefore $d_{ijk} = -d_{ijk} \implies d_{ijk} = 0$
\end{itemize}

For zinc blende (cubic but no inversion center):
\begin{itemize}
    \item The symmetry operations do not include inversion
    \item Therefore $d_{ijk}$ can be non-zero
    \item Zinc blende is indeed piezoelectric
\end{itemize}

\textbf{Grading Rubric:}
\begin{table}[h]
\centering
\begin{tabular}{|p{10cm}|c|}
\hline
\textbf{Criterion} & \textbf{Points} \\
\hline
Correct statement of Neumann's principle & 1 \\
Correct analysis of inversion symmetry & 1.5 \\
Correct analysis of piezoelectric tensor & 1.5 \\
Correct conclusion for zinc blende & 1 \\
\hline
\end{tabular}
\end{table}

\section{Total Points Summary}

\begin{table}[h]
\centering
\begin{tabular}{|c|c|c|}
\hline
\textbf{Problem} & \textbf{Topic} & \textbf{Points} \\
\hline
1 & FCC Packing Fraction & 5 \\
2 & Diamond Atoms & 3 \\
3 & (110) Plane in BCC & 4 \\
4 & Miller Indices & 4 \\
5 & Interplanar Spacing & 4 \\
6 & $(11\bar{2})$ Plane & 4 \\
7 & 2D Hexagonal Reciprocal & 8 \\
8 & BCC $\leftrightarrow$ FCC & 8 \\
9 & $\mathbf{G} \perp$ Plane & 6 \\
10 & Bragg Angle & 5 \\
11 & Diamond Structure Factor & 6 \\
12 & (100) Absence & 4 \\
13 & Powder Pattern Indexing & 6 \\
14 & Powder Method Explanation & 4 \\
15 & Multiplicity Factor & 4 \\
16 & 5-fold Rotation & 5 \\
17 & Cube Point Group & 3 \\
18 & 2D Space Groups & 2 \\
19 & Copper Properties & 9 \\
20 & BCC Diffraction & 6 \\
21 & 2D Hexagonal (Repeated) & 6 \\
22 & Sodium BCC & 6 \\
23 & Zinc Blende Structure & 7 \\
24 & Neumann's Principle & 5 \\
\hline
\multicolumn{2}{|c|}{\textbf{Total}} & \textbf{124} \\
\hline
\end{tabular}
\end{table}

\section{TA Grading Guidelines}

\subsection{General Principles}
\begin{enumerate}
    \item \textbf{Consistency:} Apply the same standards to all student submissions.
    \item \textbf{Partial Credit:} Award partial credit for correct intermediate steps.
    \item \textbf{Units:} Deduct 0.5 points for missing or incorrect units.
    \item \textbf{Significant Figures:} Deduct 0.5 points for significant figure errors.
    \item \textbf{Clarity:} Deduct up to 1 point for unclear or disorganized solutions.
\end{enumerate}

\subsection{Common Mistakes by Topic}

\subsubsection{Crystal Structures}
\begin{itemize}
    \item Mixing up FCC and BCC coordination numbers
    \item Incorrect number of atoms per unit cell
    \item Wrong packing fraction formulas
\end{itemize}

\subsubsection{Miller Indices}
\begin{itemize}
    \item Forgetting to take reciprocals
    \item Not clearing fractions properly
    \item Confusing families ($\{hkl\}$) with specific planes $(hkl)$
\end{itemize}

\subsubsection{Reciprocal Lattice}
\begin{itemize}
    \item Incorrect application of $\mathbf{a}_i \cdot \mathbf{b}_j = 2\pi\delta_{ij}$
    \item Wrong cross product calculations
    \item Forgetting the $2\pi$ factor
\end{itemize}

\subsubsection{Diffraction}
\begin{itemize}
    \item Confusing Bragg angle $\theta$ with $2\theta$
    \item Incorrect structure factor calculations
    \item Missing systematic absence conditions
\end{itemize}

\subsection{Feedback Template}
For efficient grading, use the following template for feedback:

\begin{verbatim}
Problem X:
- Approach: [Excellent/Good/Needs Improvement]
- Calculation: [Correct/Minor errors/Major errors]
- Units: [Correct/Incorrect]
- Score: [X/Y points]

Comments: [Specific feedback]
\end{verbatim}

\section{Supplementary Teaching Materials}

\subsection{Key Formulas Quick Reference}
\begin{align*}
\text{Lattice: } & \mathbf{R} = n_1\mathbf{a}_1 + n_2\mathbf{a}_2 + n_3\mathbf{a}_3 \\
\text{Reciprocal: } & \mathbf{G} = h\mathbf{b}_1 + k\mathbf{b}_2 + l\mathbf{b}_3 \\
\text{Orthogonality: } & \mathbf{a}_i \cdot \mathbf{b}_j = 2\pi\delta_{ij} \\
\text{Bragg's Law: } & 2d\sin\theta = n\lambda \\
\text{Interplanar: } & d_{hkl} = \frac{a}{\sqrt{h^2+k^2+l^2}} \\
\text{Structure Factor: } & F_{hkl} = \sum_j f_j e^{2\pi i(hx_j+ky_j+lz_j)}
\end{align*}

\subsection{Structure Factor Rules Table}
\begin{table}[h]
\centering
\begin{tabular}{|c|c|}
\hline
\textbf{Structure} & \textbf{Reflection Conditions} \\
\hline
SC & All $(hkl)$ \\
BCC & $h+k+l$ even \\
FCC & $h,k,l$ all even or all odd \\
Diamond & $h,k,l$ all even or all odd, and $h+k+l$ not $4n+2$ \\
NaCl & $h,k,l$ all even or all odd \\
CsCl & All $(hkl)$ \\
\hline
\end{tabular}
\end{table}

\end{document}